% formal/spinhint.html
% mainfile: ../perfbook.tex
% SPDX-License-Identifier: CC-BY-SA-3.0

\section{状态空间搜索}
\label{sec:formal:State-Space Search}
%
\epigraph{探索每条小径 / 走遍你知晓的每条路。}
	 {\emph{Climb Every Mountain}, Rodgers \& Hammerstein}

本节介绍通用的 Promela 和 Spin 工具，
它们可用于对多线程代码进行多种类型的完整状态空间搜索。
对于数据通信协议验证来说，它们也非常有用。
\Cref{sec:formal:Promela and Spin}
介绍了 Promela 和 Spin，包括两个热身练习——
分别验证非原子递增和原子递增。
\Cref{sec:formal:How to Use Promela}
描述了 Promela 的使用方法，包括示例命令行以及
Promela 语法与 C 语法的比较。
\Cref{sec:formal:Promela Example: Locking}
展示了如何使用 Promela 验证锁机制，
\cref{sec:formal:Promela Example: QRCU}
使用 Promela 验证一种名为``QRCU''的特殊 RCU 实现，
最后
\cref{sec:formal:Promela Parable: dynticks and Preemptible RCU}
将 Promela 应用于 RCU dyntick-idle 实现的早期版本。

\subsection{Promela 和 Spin}
\label{sec:formal:Promela and Spin}

\IX{Promela} 是一种旨在帮助验证协议的语言，但也可以用于验证小型并行算法。
你需要用类似于 C 语言的 Promela 来重新编写你的算法和正确性约束，
然后使用 \IX{Spin} 将其翻译为可以编译和运行的 C 程序。
生成的程序会对你的算法执行完整的状态空间搜索，
要么确认通过你在 Promela 程序中关联的断言，要么找到反例。

这种完整的状态搜索非常强大，但也是一把双刃剑。
如果算法太复杂，或者 Promela 实现得不够仔细，
那么状态空间可能大到内存无法容纳。
此外，即使有充足的内存，状态空间搜索的运行时间
也可能超过宇宙的预期寿命。
因此，请在复杂但紧凑的并行算法中使用此工具。
轻率地试图将其应用于中等规模的算法（更不用说整个 Linux 内核）
将会以失败告终。

Promela 和 Spin 可以从
\url{https://spinroot.com/spin/whatispin.html} 下载。

上述网站还提供了 Gerard Holzmann 关于 Promela 和 Spin 的优秀
著作~\cite{Holzmann03a}的链接，
可以通过以下地址进行在线搜索：
\url{https://www.spinroot.com/spin/Man/index.html}。

本节其余部分描述如何使用 Promela 调试并行算法，
从简单示例开始，然后逐步过渡到更复杂的使用场景。

\subsubsection{热身：
			非原子递增}
\label{sec:formal:Warm-Up: Non-Atomic Increment}

\begin{fcvref}[ln:formal:promela:increment:whole]
\Cref{lst:formal:Promela Code for Non-Atomic Increment}
演示了教科书中所讲的非原子递增导致的\IX{竞态条件}。
\Clnref{nprocs} 定义了要运行的进程数量（我们将改变此值
以观察对状态空间的影响），\clnref{count} 定义了计数器，
\clnref{prog} 用于实现
\clnrefrange{assert:b}{assert:e} 上的断言。

\begin{listing}
\input{CodeSamples/formal/promela/increment@whole.fcv}
\caption{非原子递增的 Promela 代码}
\label{lst:formal:Promela Code for Non-Atomic Increment}
\end{listing}

\Clnrefrange{proc:b}{proc:e} 定义了一个非原子地递增计数器的过程。
参数 \co{me} 是进程编号，由后面的初始化代码进行设置。
由于简单的 Promela 语句各自被假定为原子的，
我们必须将递增操作拆分为
\clnrefrange{incr:b}{incr:e} 上的两条语句。
\clnref{setprog} 上的赋值标记该过程的完成。
由于 Spin 系统会搜索整个状态空间，包括所有可能的状态序列，
因此不需要传统压力测试中使用的循环。

\Clnrefrange{init:b}{init:e} 是初始化代码块，它最先被执行。
\Clnrefrange{doinit:b}{doinit:e} 实际进行初始化操作，
而 \clnrefrange{assert:b}{assert:e} 执行断言。
两者都是原子块，以避免不必要地增加状态空间：
因为它们不属于算法本身的一部分，
将它们标记为原子的不会导致验证覆盖率的缺失。

\clnrefrange{dood1:b}{dood1:e} 上的 \co{do-od} 构造实现了一个 Promela 循环，
可以将其理解为一个 C 语言的 \co{for (;;)} 循环，
其中包含一个允许在 case 标签中使用表达式的 \co{switch} 语句。
条件块（以 \co{::} 为前缀）会被非确定性地扫描，
不过在本例中，在任何给定时刻只有一个条件可能成立。
\co{do-od} 的第一个块从
\clnrefrange{block1:b}{block1:e}，
初始化第 i 个递增器的进度单元，运行第 i 个递增器的进程，
然后递增变量 \co{i}。
\co{do-od} 的第二个块在 \clnref{block2} 上，
在这些进程启动后退出循环。

\clnrefrange{assert:b}{assert:e} 上的原子块也包含
一个类似的 \co{do-od} 循环，用于累加进度计数器。
\clnref{assert} 上的 \co{assert()} 语句验证：
如果所有进程都已完成，则所有计数都已正确记录。
\end{fcvref}

你可以按如下方式构建并运行此程序：

\begin{VerbatimU}
spin -a increment.spin      # Translate the model to C
cc -DSAFETY -o pan pan.c    # Compile the model
./pan                       # Run the model
\end{VerbatimU}

\begin{listing}
\VerbatimInput[numbers=none,fontsize=\scriptsize]{CodeSamples/formal/promela/increment.spin.lst}
\vspace*{-9pt}
\caption{非原子递增的 Spin 输出}
\label{lst:formal:Non-Atomic Increment Spin Output}
\end{listing}

这将产生如
\cref{lst:formal:Non-Atomic Increment Spin Output}
所示的输出。
第一行告诉我们断言被违反了（考虑到非原子递增，这在预期之中！）。
第二行中的跟踪（trail）文件描述了断言是如何被违反的。
``Warning''行重申在我们的模型中，并非所有事情都正确完成。
第二段描述了执行的状态搜索类型，
本例中是检查断言违反和无效终止状态。
第三段给出了状态大小统计信息：
这个小模型只有 45 个状态。
最后一行显示了内存使用情况。

跟踪（trail）文件可以被整理为如下所示的人类可读格式：

\begin{VerbatimU}
spin -t -p increment.spin
\end{VerbatimU}

\begin{listing*}
\ebresizeverb{.9}{
\VerbatimInput[numbers=none,fontsize=\scriptsize]{CodeSamples/formal/promela/increment.spin.trail.lst}}
\IfEbookSize{\vspace*{7pt}}{\vspace*{-9pt}}
\caption{非原子递增错误轨迹}
\label{lst:formal:Non-Atomic Increment Error Trail}
\end{listing*}

它给出了如
\cref{lst:formal:Non-Atomic Increment Error Trail}
所示的输出。
正如我们所见，初始化块的第一部分创建了两个递增器进程，
其中第一个进程获得了计数器的值，然后两个进程都递增并存储了它，
从而丢失了一次计数。
随后断言被触发，之后显示了全局状态。

\subsubsection{热身：
			原子递增}
\label{sec:formal:Warm-Up: Atomic Increment}

像\cref{lst:formal:Promela Code for Atomic Increment}那样，
通过将递增过程的主体放入原子块中，可以简单地修复此示例。
由于 Promela 语句是原子的，一个简单的方法是用
\co{counter = counter + 1} 语句来替换递增语句对。
无论哪种方式，运行修改后的模型都会给出无错误的状态空间遍历，
如\cref{lst:formal:Atomic Increment Spin Output}所示。

\begin{listing}
\input{CodeSamples/formal/promela/atomicincrement@incrementer.fcv}
\caption{原子递增的 Promela 代码}
\label{lst:formal:Promela Code for Atomic Increment}
\end{listing}

\begin{listing}
\VerbatimInput[numbers=none,fontsize=\scriptsize]{CodeSamples/formal/promela/atomicincrement.spin.lst}
\vspace*{-9pt}
\caption{原子递增的 Spin 输出}
\label{lst:formal:Atomic Increment Spin Output}
\end{listing}

\Cref{tab:advsync:Memory Usage of Increment Model}
显示了作为所建模的递增器数量（通过重新定义 \co{NUMPROCS}）
的函数的状态数和内存消耗：

\begin{table}
\rowcolors{1}{}{lightgray}
\small
\renewcommand*{\arraystretch}{1.2}
\centering
\begin{tabular}{S[table-format = 1.0]S[table-format = 7.0]S[table-format = 3.1]}
	\toprule
	\multicolumn{1}{l}{递增器数量} &
		\multicolumn{1}{r}{状态数} &
			\multicolumn{1}{r}{总内存使用量 (MB)} \\
	\midrule
	1 &		        11 &        128.7 \\
	2 &		        52 &        128.7 \\
	3 &		       372 &        128.7 \\
	4 &		     3 496 &        128.9 \\
	5 &		    40 221 &        131.7 \\
	6 &		   545 720 &        174.0 \\
	7 &		 8 521 446 &        881.9 \\
	\bottomrule
\end{tabular}
\caption{递增模型的内存使用量}
\label{tab:advsync:Memory Usage of Increment Model}
\end{table}

因此，不必要地运行大型模型的前景看起来不容乐观，尽管
882\,MB 完全在现代台式机和笔记本电脑的限制范围内。

有了这个示例的经验，让我们更仔细地看看
用于分析 Promela 模型的命令，然后看看更精细的示例。

\subsection{如何使用 Promela}
\label{sec:formal:How to Use Promela}

给定一个源文件 \path{qrcu.spin}，可以使用以下命令：

\begin{description}[style=nextline]
\item	[\tco{spin -a qrcu.spin}]
	创建一个完整搜索状态机的文件 \path{pan.c}。
\item	[\tco{cc -DSAFETY [-DCOLLAPSE] [-DMA=N] -o pan pan.c}]
	编译生成的状态机搜索程序。
	\co{-DSAFETY} 生成适用于仅包含断言（以及可能的 \co{never} 语句）
	情况下的优化。
	如果你有活性、公平性或前向进展检查，
	可能需要在不使用 \co{-DSAFETY} 的情况下编译。
	如果你在可以使用 \co{-DSAFETY} 时没有使用它，
	程序会提示你。

	\co{-DSAFETY} 产生的优化大大加快了速度，
	所以应尽可能使用它。
	不能使用 \co{-DSAFETY} 的一个示例场景是
	通过 \co{-DNP} 检查\IXpl{livelock}（又称``非进展循环''）时。

	可选的 \co{-DCOLLAPSE} 标志生成状态向量压缩模式的代码。

	另一个可选标志 \co{-DMA=N} 生成一种较慢但更激进的
	状态空间内存压缩模式的代码。
\item	[\tco{./pan [-mN] [-wN]}]
	这将实际搜索状态空间。
	即使对于非常小的状态机，状态数也可以达到数千万，
	因此你需要一台具有大内存的机器。
	例如，具有 3 个更新者和 2 个读者的 \path{qrcu.spin}
	即使使用 \co{-DCOLLAPSE} 标志也需要 10.5\,GB 内存。

	如果你看到 \co{./pan} 显示的消息：
	``\co{error: max search depth too small}''，你需要通过
	\co{-mN} 选项增加最大深度以进行完整搜索。
	默认值为 \co{-m10000}。

	\co{-wN} 选项指定哈希表（hash table）大小。
	完整状态空间搜索的默认值为 \co{-w24}。\footnote{
		截至 Spin 版本 6.4.6 和 6.4.8。
		在 2011 年 7 月 10 日的 Spin 在线手册中，
		穷举搜索模式的默认值据说是 \co{-w19}，
		这与实际行为不符。}

	如果你不确定你的机器是否有足够的内存，
	在一个窗口中运行 \co{top}，在另一个窗口中运行 \co{./pan}。
	将焦点保持在 \co{./pan} 窗口上，以便在需要时快速终止执行。
	一旦 CPU 时间明显低于 100\pct，就终止 \co{./pan}。
	如果你已将焦点从运行 \co{./pan} 的窗口移开，
	可能需要等待很长时间，窗口系统才能获得足够的内存来为你做任何事情。

	另一种避免内存耗尽的选择是
	\co{-DMEMLIM=N} 编译器标志。
	\co{-DMEMLIM=2000} 将设置 2\,GB 的最大值。

	不要忘记捕获输出，特别是在远程机器上工作时。

	如果你的模型包含前向进展检查，你可能需要
	通过 \co{./pan} 的 \co{-f} 命令行参数启用``弱公平性''。
	如果你的前向进展检查涉及 \co{accept} 标签，
	你还需要 \co{-a} 参数。
	% forward reference to model: formal.2009.02.19a in
	% /home/linux/git/userspace-rcu/formal-model.
\item	[\tco{spin -t -p qrcu.spin}]
	给定遇到错误的运行所输出的 \co{trail} 文件，
	输出导致该错误的步骤序列。
	\co{-g} 标志还会包含已更改的全局变量的值，
	\co{-l} 标志还会包含已更改的局部变量的值。
\end{description}

\subsubsection{Promela 的特殊之处}
\label{sec:formal:Promela Peculiarities}

尽管所有计算机语言都有潜在的相似性，
但 Promela 仍然会给习惯于用 C、C++ 或 Java 编程的人带来一些意外。

\begin{enumerate}
\item	在 C 中，\qco{;} 终止一条语句。
	在 Promela 中，它将语句进行分隔。
	幸运的是，较新版本的 Spin 已经对此比较宽容了。
\item	Promela 的循环构造 \co{do} 语句需要循环条件。
	该 \co{do} 语句类似于 if-then-else 语句的循环版本。
\item	在 C 的 \co{switch} 语句中，如果没有匹配的情况，
	整个语句被忽略。
	在 Promela 中，其等价物令人困惑地被称为 \co{if}，
	如果没有相应的守卫表达式匹配，将得到一个错误，
	而且没有相应的可识别的错误消息。
	因此，如果错误输出指向无辜的代码行，
	请检查你是否在 \co{if} 或 \co{do} 语句中遗漏了某个条件。
\item	在 C 中创建压力测试时，通常会将可疑操作相互放在一起运行。
	在 Promela 中，取而代之的是只运行一次操作，因为 Promela
	会搜索出该操作所有可能的结果。
	有时你确实需要在 Promela 中循环操作，例如，
	当多个操作交叉时，但这样做会大大增加状态空间的大小。
\item	在 C 中，最简单的做法是维护一个循环计数器来跟踪进度和终止循环。
	在 Promela 中，循环计数器必须像躲避瘟疫一样避免使用，
	因为它们会导致状态空间迅速膨胀。
	但从另一方面来说，在 Promela 中使用无限循环是没有问题的，
	只要循环里面没有变量单调递增或递减——Promela 会自行计算出
	循环到底需要多少次遍历才有意义，
	并自动删除超出该点的不必要执行。
\item	在 C 压力测试代码中，使用每任务控制变量通常是明智的。
	它们读取开销低，对调试测试代码也很有帮助。
	在 Promela 中，仅仅在没有其他替代方案时，才使用每任务变量。
	要明白这一点，考虑一个 5 任务的验证，
	每一位表示一个任务的完成。
	这会导致 32 个状态。
	相比之下，一个简单的计数器只有六个状态，
	状态数整整少了五倍多。
	五倍看起来可能不是什么问题，但如果你在为一个拥有超过
	1.5 亿个状态、消耗超过 10\,GB 内存的验证程序而挣扎时，
	就不会这样想了！
\item	无论是在 C 压力测试代码还是 Promela 中，最具挑战性的事情之一
	是编写好的断言。
	Promela 还允许 \co{never} 声明，它的作用类似于
	在每一行代码之间复制的断言。
\item	在 Promela 中，分而治之对于将状态空间保持在可控范围内非常有用。
	将一个大型模型分成大致相等的两半，
	会将状态空间减少到大约整体的平方根。
	例如，一个百万状态的组合模型可能缩减为
	一对千状态的模型。
	Promela 不仅能以更少的内存更快地处理两个较小的模型，
	而且两个较小的算法也更容易让人理解。
\end{enumerate}


\subsubsection{Promela 编码技巧}
\label{sec:formal:Promela Coding Tricks}

Promela 被设计用于分析协议，因此将其用于并行程序算是有点滥用。
以下技巧可以帮助你安全地用好 Promela：

\begin{enumerate}
\item	内存乱序。
	假设你有一对语句将全局变量 \co{x} 和 \co{y}
	复制到局部变量 \co{r1} 和 \co{r2}，其中顺序很重要（例如，没有锁保护），
	但你没有\IXpl{memory barrier}。
	这可以在 Promela 中用如下方法建模：

\begin{VerbatimN}[samepage=true]
if
:: 1 -> r1 = x;
        r2 = y
:: 1 -> r2 = y;
        r1 = x
fi
\end{VerbatimN}

	\co{if} 语句的两个分支会被非确定性地选取，因为它们都是可用的。
	由于进行了完整的状态空间搜索，在所有情况下
	\emph{两种}选择最终都会被覆盖。

	当然，过度使用这个技巧会导致你的状态空间迅速增大。
	此外，它要求你预料到可能的乱序情况。

\item	状态压缩。
	如果你有复杂的断言，请在 \co{atomic} 块内求值。
	毕竟，它们不是算法的一部分。
	一个复杂断言的例子（后面将更详细地讨论）
	如\cref{lst:formal:Complex Promela Assertion}所示。

	没有理由非原子地执行这个断言，
	因为它实际上不是算法的一部分。
	因为每条语句都会增加状态数量，我们可以通过将其
	放入 \co{atomic} 块中来减少无用状态的数量，
	如\cref{lst:formal:Atomic Block for Complex Promela Assertion}所示。

\item	Promela 不提供函数。
	必须使用 C 预处理器宏来代替。
	但是，要小心使用它们以避免状态组合快速增加。
\end{enumerate}

\begin{listing}
\begin{VerbatimL}
i = 0;
sum = 0;
do
:: i < N_QRCU_READERS ->
	sum = sum + (readerstart[i] == 1 &&
	             readerprogress[i] == 1);
	i++
:: i >= N_QRCU_READERS ->
	assert(sum == 0);
	break
od
\end{VerbatimL}
\caption{复杂的 Promela 断言}
\label{lst:formal:Complex Promela Assertion}
\end{listing}

\begin{listing}
\begin{VerbatimL}
atomic {
	i = 0;
	sum = 0;
	do
	:: i < N_QRCU_READERS ->
		sum = sum + (readerstart[i] == 1 &&
		             readerprogress[i] == 1);
		i++
	:: i >= N_QRCU_READERS ->
		assert(sum == 0);
		break
	od
}
\end{VerbatimL}
\caption{复杂 Promela 断言的原子块}
\label{lst:formal:Atomic Block for Complex Promela Assertion}
\end{listing}

现在我们准备好见识更复杂的示例了。

\subsection{Promela 示例：
			     锁}
\label{sec:formal:Promela Example: Locking}

\begin{fcvref}[ln:formal:promela:lock:whole]
由于锁通常是有用的，\path{lock.h} 中提供了 \co{spin_lock()} 和
\co{spin_unlock()} 宏，可以从多个 Promela 模型中包含，
如\cref{lst:formal:Promela Code for Spinlock}所示。
\co{spin_lock()} 宏包含一个无限 \co{do-od} 循环，
跨越 \clnrefrange{dood:b}{dood:e}，
由 \clnref{one} 上的单一守卫表达式``1''驱动。
该循环的主体是一个包含 \co{if-fi} 语句的原子块。
\co{if-fi} 构造与 \co{do-od} 构造类似，
只是它只执行一次而不是循环。
如果锁在 \clnref{notheld} 上未被持有，则
\clnref{acq} 获取它，
\clnref{break} 跳出外层 \co{do-od} 循环（同时也退出原子块）。
另一方面，如果锁在 \clnref{held} 上已被持有，
我们什么都不做（\co{skip}），并跳出 \co{if-fi} 和原子块，
以便再次通过外层循环，重复直到锁可用。
\end{fcvref}

\begin{listing}
\input{CodeSamples/formal/promela/lock@whole.fcv}
\caption{自旋锁（spinlock）的 Promela 代码}
\label{lst:formal:Promela Code for Spinlock}
\end{listing}

\co{spin_unlock()} 宏简单地将锁标记为不再被持有。

请注意，不需要内存屏障（memory barrier），因为 Promela 假设强序。
在任何给定的 Promela 状态中，所有处理器对当前状态以及
导致我们到达当前状态的状态变化顺序都达成一致。
这类似于少数计算机系统（如 1990 年代的 MIPS 和 PA-RISC）
所使用的``\IXalth{sequentially consistent}{sequential}{memory consistency}''
\IX{memory model}。
如前面的提示所述，并且将在后面的示例中看到，
弱内存排序必须被显式编码。

\begin{listing}
\input{CodeSamples/formal/promela/lock@spin.fcv}
\caption{测试自旋锁的 Promela 代码}
\label{lst:formal:Promela Code to Test Spinlocks}
\end{listing}

\begin{fcvref}[ln:formal:promela:lock:spin]
这些宏通过\cref{lst:formal:Promela Code to Test Spinlocks}
中所示的 Promela 代码进行测试。
此代码与用于测试递增的代码类似，
锁定进程的数量由 \clnref{nlockers} 上的
\co{N_LOCKERS} 宏定义来定义。
互斥锁（exclusive lock）本身在 \clnref{mutex} 上定义，
用于跟踪锁持有者的数组在 \clnref{array} 上，
\clnref{sum} 由断言代码使用以验证只有一个进程持有锁。
\end{fcvref}

\begin{fcvref}[ln:formal:promela:lock:spin:locker]
锁定器进程在 \clnrefrange{b}{e} 上，简单地无限循环：
在 \clnref{lock} 上获取锁，在 \clnref{claim} 上声明持有，
在 \clnref{unclaim} 上取消声明，在 \clnref{unlock} 上释放锁。
\end{fcvref}

\begin{fcvref}[ln:formal:promela:lock:spin:init]
\clnrefrange{b}{e} 上的 init 块在 \clnref{array} 上初始化当前锁定器的
\co{havelock} 数组项，在 \clnref{start} 上启动当前锁定器，
并在 \clnref{next} 上前进到下一个锁定器。
一旦所有锁定器进程被创建，\co{do-od} 循环
移动到 \clnref{chkassert}，用于检查断言。
\Clnref{sum,j} 初始化控制变量，
\clnrefrange{atm:b}{atm:e} 原子地对 \co{havelock} 数组项求和，
\clnref{assert} 是断言，\clnref{break} 退出循环。
\end{fcvref}

我们可以通过将
\cref{lst:formal:Promela Code for Spinlock,%
lst:formal:Promela Code to Test Spinlocks} 中的两个代码片段
分别放入名为 \path{lock.h} 和 \path{lock.spin} 的文件中，
然后运行以下命令来运行此模型：

\begin{VerbatimU}
spin -a lock.spin
cc -DSAFETY -o pan pan.c
./pan
\end{VerbatimU}

\begin{listing}
\VerbatimInput[numbers=none,fontsize=\scriptsize]{CodeSamples/formal/promela/lock.spin.lst}
\vspace*{-9pt}
\caption{自旋锁测试的输出}
\label{lst:formal:Output for Spinlock Test}
\end{listing}

输出将看起来类似于
\cref{lst:formal:Output for Spinlock Test}所示。
如预期的那样，此运行没有断言失败（\qco{errors: 0}）。

\QuickQuizSeries{%
\QuickQuizB{
	为什么 locker 中有一条未到达的语句？
	毕竟，这不是\emph{完整的}状态空间搜索吗？
}\QuickQuizAnswerB{
	locker 进程是一个无限循环，因此控制永远不会
	到达该进程的末尾。
	但是，由于没有单调递增的变量，
	Promela 能够用少量状态来建模这个无限循环。
}\QuickQuizEndB
%
\QuickQuizE{
	这个示例有哪些 Promela 代码风格问题？
}\QuickQuizAnswerE{
	有几个：
	\begin{enumerate}
	\item	\co{sum} 的声明应该移到 init 块内部，
		因为它在其他地方没有被使用。
	\item	断言代码应该移到初始化循环之外。
		然后可以将初始化循环放在原子块中，
		大大减少状态空间（减少多少？）。
	\item	覆盖断言代码的原子块应该扩展到包含
		\co{sum} 和 \co{j} 的初始化，
		并且也覆盖断言。
		这也减少了状态空间（又减少了多少？）。
	\end{enumerate}
}\QuickQuizEndE
}

\subsection{Promela 示例：
			     QRCU}
\label{sec:formal:Promela Example: QRCU}

这个最终示例演示了在 Oleg Nesterov 的
QRCU~\cite{OlegNesterov2006QRCU,OlegNesterov2006aQRCU}
上实际使用 Promela 的情况，
但进行了修改以加速 \co{synchronize_qrcu()} 的快速路径。

但首先需要搞清楚，QRCU 是什么？

QRCU 是 SRCU~\cite{PaulEMcKenney2006c} 的一个变体，
它以较高的读开销（对全局变量的原子递增和递减）换取
极低的宽限期延迟。
如果没有读者，宽限期将在不到一微秒内被检测到，
而大多数其他 RCU 实现的宽限期延迟则为数毫秒。

\begin{enumerate}
\item	有一个 \co{qrcu_struct} 定义了一个 QRCU 域。
	与 SRCU 类似（不同于其他 RCU 变体），QRCU 的作用
	不是全局的，而是集中在指定的 \co{qrcu_struct} 上。
\item	有 \co{qrcu_read_lock()} 和 \co{qrcu_read_unlock()}
	原语用于界定 QRCU 读端临界区（critical section）。
	必须将相应的 \co{qrcu_struct} 传递给这些原语，
	并且 \co{qrcu_read_lock()} 的返回值
	必须传递给 \co{qrcu_read_unlock()}。

	例如：

\begin{VerbatimU}
idx = qrcu_read_lock(&my_qrcu_struct);
/* read-side critical section. */
qrcu_read_unlock(&my_qrcu_struct, idx);
\end{VerbatimU}

\item	有一个 \co{synchronize_qrcu()} 原语，它会阻塞直到
	所有先前存在的 QRCU 读端临界区完成，
	但与 SRCU 的 \co{synchronize_srcu()} 类似，QRCU 的
	\co{synchronize_qrcu()} 只需等待使用相同 \co{qrcu_struct}
	的那些读端临界区。

	例如，\co{synchronize_qrcu(&your_qrcu_struct)}
	\emph{不}需要等待前面的 QRCU 读端临界区。
	相反，\co{synchronize_qrcu(&my_qrcu_struct)}
	\emph{需要}等待，因为它共享相同的 \co{qrcu_struct}。
\end{enumerate}

已经产生了一个 QRCU 的 Linux 内核补丁~\cite{PaulMcKenney2007QRCUpatch}，
但它不太可能被包含在 Linux 内核中。

\begin{listing}
\input{CodeSamples/formal/promela/qrcu@gvar.fcv}
\caption{QRCU 全局变量}
\label{lst:formal:QRCU Global Variables}
\end{listing}

回到 QRCU 的 Promela 代码，全局变量如
\cref{lst:formal:QRCU Global Variables}所示。
这个示例使用了锁并包含了 \path{lock.h}。
读者和写者的数量都可以通过两个 \co{#define} 语句来改变，
这给了我们不止一种而是两种造成组合爆炸的方式。
\co{idx} 变量控制 \co{ctr} 数组的两个元素中哪一个将被读者使用，
\co{readerprogress} 变量允许断言判断所有读者何时完成
（因为直到每个已存在的读者完成其 QRCU 读端临界区之前，
QRCU 更新都不允许完成）。
\co{readerprogress} 数组元素的值如下，
表示相应读者的状态：

\begin{enumerate}[label={\arabic*}:,start=0,itemsep=0pt]
\item	尚未开始。
\item	在 QRCU 读端临界区内。
\item	已完成 QRCU 读端临界区。
\end{enumerate}

最后，\co{mutex} 变量用于序列化更新者的慢速路径。

\begin{listing}
\input{CodeSamples/formal/promela/qrcu@reader.fcv}
\caption{QRCU 读者进程}
\label{lst:formal:QRCU Reader Process}
\end{listing}

\begin{fcvref}[ln:formal:promela:qrcu:reader]
QRCU 读者由\cref{lst:formal:QRCU Reader Process}中所示的
\co{qrcu_reader()} 进程建模。
一个 \co{do-od} 循环跨越 \clnrefrange{do}{od}，
在 \clnref{one} 上有一个值为``1''的单一守卫，
使其成为无限循环。
\Clnref{curidx} 捕获全局索引的当前值，
\clnrefrange{atm:b}{atm:e} 在其值非零时原子地递增它
（并跳出无限循环）（\co{atomic_inc_not_zero()}）。
\Clnref{cs:entry} 标记进入 RCU 读端临界区，
\clnref{cs:exit} 标记退出该临界区，
这两行都是为了后面将遇到的 \co{assert()} 语句服务的。
\Clnref{atm:dec} 原子地递减我们之前递增的同一计数器，
从而退出 RCU 读端临界区。
\end{fcvref}

\begin{listing}
\input{CodeSamples/formal/promela/qrcu@sum_unordered.fcv}
\caption{QRCU 无序求和}
\label{lst:formal:QRCU Unordered Summation}
\end{listing}

\begin{fcvref}[ln:formal:promela:qrcu:sum_unordered]
\cref{lst:formal:QRCU Unordered Summation}
中所示的 C 预处理器宏对一对计数器求和，
以模拟弱内存排序。
\Clnrefrange{fetch:b}{fetch:e} 获取其中一个计数器，
\clnref{sum_other} 获取另一个并求和。
原子块由单个 \co{do-od} 语句组成。
这个 \co{do-od} 语句（跨越 \clnrefrange{do}{od}）不寻常之处在于
它包含两个在 \clnref{g1,g2} 上带有守卫的无条件分支，
这导致 Promela 非确定性地选择其中一个（但同样，
完整的状态空间搜索使 Promela 最终在每个适用情况下做出所有可能的选择）。
第一个分支获取第零个计数器并将 \co{i} 设置为 1（以便
\clnref{sum_other} 获取第一个计数器），而第二个分支相反，
获取第一个计数器并将 \co{i} 设置为 0（以便
\clnref{sum_other} 获取第二个计数器）。
\end{fcvref}

\QuickQuiz{
	有没有更直接的方式来编写 \co{do-od} 语句？
}\QuickQuizAnswer{
	有。
	将其替换为 \co{if-fi} 并删除两个 \co{break} 语句。
}\QuickQuizEnd

\begin{listing}
\ebresizeverb{.88}{\input{CodeSamples/formal/promela/qrcu@updater.fcv}}
\caption{QRCU 更新者进程}
\label{lst:formal:QRCU Updater Process}
\end{listing}

\begin{fcvref}[ln:formal:promela:qrcu:updater]
有了 \co{sum_unordered} 宏，我们现在可以继续看
\cref{lst:formal:QRCU Updater Process}中所示的更新端进程。
更新端进程无限重复，相应的 \co{do-od} 循环
跨越 \clnrefrange{do}{od}。
每次循环首先在 \clnrefrange{atm1:b}{atm1:e} 上
将全局 \co{readerprogress} 数组快照到局部 \co{readerstart} 数组中。
此快照将用于 \clnref{assert} 上的断言。
\Clnref{sum_unord} 调用 \co{sum_unordered}，然后
\clnrefrange{reinvoke:b}{reinvoke:e}
在快速路径可能可用时重新调用 \co{sum_unordered}。

\Clnrefrange{slow:b}{slow:e} 在需要时执行慢速路径代码，
\clnref{acq,rel} 获取和释放更新端锁，
\clnrefrange{flip_idx:b}{flip_idx:e} 翻转索引，
\clnrefrange{wait:b}{wait:e} 等待所有先前存在的读者完成。

\Clnrefrange{atm2:b}{atm2:e} 然后将 \co{readerprogress}
数组中的当前值与 \co{readerstart} 数组中收集的值进行比较，
如果在此更新之前开始的任何读者仍在进行中，
则强制断言失败。
\end{fcvref}

\QuickQuizSeries{%
\QuickQuizB{
	\begin{fcvref}[ln:formal:promela:qrcu:updater]
	为什么在 \clnrefrange{atm1:b}{atm1:e}
	和 \clnrefrange{atm2:b}{atm2:e} 上有原子块，
	而这些原子块中的操作在当前任何生产级微处理器上
	都没有原子实现？
	\end{fcvref}
}\QuickQuizAnswerB{
	因为这些操作仅用于断言。
	它们不是算法本身的一部分。
	因此将它们标记为原子的没有害处，
	这样标记大大减少了 Promela 模型必须搜索的状态空间。
}\QuickQuizEndB
%
\QuickQuizE{
	\begin{fcvref}[ln:formal:promela:qrcu:updater]
	\clnrefrange{reinvoke:b}{reinvoke:e}
	上的计数器重新求和\emph{真的}必要吗？
	\end{fcvref}
}\QuickQuizAnswerE{
	是的。
	要理解这一点，请删除这些行并运行模型。

	或者，考虑以下步骤序列：

	\begin{enumerate}
	\item	一个进程在其 RCU 读端临界区内，
		因此 \co{ctr[0]} 的值为零，
		\co{ctr[1]} 的值为二。
	\item	一个更新者开始执行，看到计数器之和为二，
		因此快速路径不能执行。
		它因此获取了锁。
	\item	第二个更新者开始执行，获取了 \co{ctr[0]} 的值，
		为零。
	\item	第一个更新者将 \co{ctr[0]} 加一，翻转索引
		（现在变为零），然后从 \co{ctr[1]} 减一
		（现在变为一）。
	\item	第二个更新者获取 \co{ctr[1]} 的值，
		现在为一。
	\item	第二个更新者现在错误地得出结论认为
		可以安全地走快速路径，尽管
		原始读者尚未完成。
	\end{enumerate}
}\QuickQuizEndE
}

\begin{listing}
\input{CodeSamples/formal/promela/qrcu@init.fcv}
\caption{QRCU 初始化进程}
\label{lst:formal:QRCU Initialization Process}
\end{listing}

\begin{fcvref}[ln:formal:promela:qrcu:init]
剩下的就是\cref{lst:formal:QRCU Initialization Process}中所示的
初始化块。
这个块只是在 \clnrefrange{i_ctr:b}{i_ctr:e} 上初始化计数器对，
在 \clnrefrange{spn_r:b}{spn_r:e} 上生成读者进程，
在 \clnrefrange{spn_u:b}{spn_u:e} 上生成更新者进程。
这些都在一个原子块内完成以减少状态空间。
\end{fcvref}

\subsubsection{运行 QRCU 示例}
\label{sec:formal:Running the QRCU Example}

要运行 QRCU 示例，请将上一节中的代码片段合并到
一个名为 \path{qrcu.spin} 的文件中，并将
\co{spin_lock()} 和 \co{spin_unlock()} 的定义放入
名为 \path{lock.h} 的文件中。
然后使用以下命令构建并运行 QRCU 模型：

\begin{VerbatimU}
spin -a qrcu.spin
cc -DSAFETY [-DCOLLAPSE] -o pan pan.c
./pan [-mN]
\end{VerbatimU}

\begin{table}
\centering
\begin{threeparttable}
\rowcolors{1}{}{lightgray}
\renewcommand*{\arraystretch}{1.2}
\footnotesize
\begin{tabular}{S[table-format = 1.0]S[table-format = 1.0]S[table-format = 9.0]
		S[table-format = 6.0]S[table-format = 5.1]}
	\toprule
	\multicolumn{1}{r}{更新者} &
	    \multicolumn{1}{r}{读者} &
		\multicolumn{1}{r}{状态数} &
		    \multicolumn{1}{r}{深度} &
			\multicolumn{1}{r}{内存 (MB)\tnote{a}} \\
	\midrule
	1 & 1 &         376 &      95 &    128.7 \\
	1 & 2 &       6 177 &     218 &    128.9 \\
	1 & 3 &      99 728 &     385 &    132.6 \\
	2 & 1 &      29 399 &     859 &    129.8 \\
	2 & 2 &   1 071 181 &   2 352 &    169.6 \\
	2 & 3 &  33 866 736 &  12 857 &  1 540.8 \\
	3 & 1 &   2 749 453 &  53 809 &    236.6 \\
	3 & 2 & 186 202 860 & 328 014 & 10 483.7 \\
	\bottomrule
\end{tabular}
\begin{tablenotes}
	\item [a] 使用编译器标志 \co{-DCOLLAPSE} 获得。
\end{tablenotes}
\end{threeparttable}
\caption{QRCU 模型的内存使用量}
\label{tab:advsync:Memory Usage of QRCU Model}
\end{table}

输出显示此模型通过了
\cref{tab:advsync:Memory Usage of QRCU Model}中所示的所有案例。
如果能运行三个读者和三个更新者就好了，
然而，简单外推表明这将需要大约半太字节的内存。
怎么办？

事实证明，当 \co{./pan} 耗尽内存时会给出建议，
例如，当试图运行三个读者和三个更新者时：

\begin{VerbatimU}
hint: to reduce memory, recompile with
  -DCOLLAPSE # good, fast compression, or
  -DMA=96   # better/slower compression, or
  -DHC # hash-compaction, approximation
  -DBITSTATE # supertrace, approximation
\end{VerbatimU}

让我们尝试建议的编译器标志 \co{-DMA=N}，
它生成以大大增加搜索开销为代价来激进压缩状态空间的代码。
所需命令如下：

\begin{VerbatimU}
spin -a qrcu.spin
cc -DSAFETY -DMA=96 -O2 -o pan pan.c
./pan -m20000000
\end{VerbatimU}

这里，深度限制 20,000,000 比从简单外推推导出的预期深度
大一个数量级。
虽然这增加了前期内存使用，但避免了由于过紧的
深度限制导致的不完整搜索而浪费漫长的运行。
此运行在一台 \Power{9} 服务器上花费了略多于 3 天。
结果如
\cref{lst:formal:spinhint:3 Readers 3 Updaters QRCU Spin Output with -DMA=96}所示。
此 Spin 运行成功完成，总内存使用量仅为 6.5\,GB，
比 \co{-DCOLLAPSE} 使用的约半太字节低了近两个数量级。

\begin{listing}
\VerbatimInput[numbers=none,fontsize=\scriptsize]{CodeSamples/formal/promela/qrcu.spin.33ma.lst}
\vspace*{-9pt}
\caption{使用 \co{-DMA=96} 的 3 读者 3 更新者 QRCU Spin 输出}
\label{lst:formal:spinhint:3 Readers 3 Updaters QRCU Spin Output with -DMA=96}
\end{listing}

\QuickQuiz{
	0.48\pct\ 的压缩率对应于状态占用内存
	200 比 1 的减少！
	状态空间搜索\emph{真的}是穷举的吗？？？
}\QuickQuizAnswer{
	根据 Spin 的文档，是的，它是穷举的。

\begin{listing}
\VerbatimInput[numbers=none,fontsize=\scriptsize]{CodeSamples/formal/promela/qrcu.spin.col-ma.diff.lst}
\vspace*{-9pt}
\caption{\co{-DCOLLAPSE} 和 \co{-DMA=88} 的 Spin 输出差异}
\label{lst:formal:promela:Spin Output Diff of -DCOLLAPSE and -DMA=88}
\end{listing}

	作为间接证据，让我们比较
	使用 \co{-DCOLLAPSE} 和 \co{-DMA=88}
	（两个读者和三个更新者）运行的结果。
	这些运行的输出差异如
	\cref{lst:formal:promela:Spin Output Diff of -DCOLLAPSE and -DMA=88}所示。
	如你所见，它们在状态数（已存储和已匹配）上一致。
}\QuickQuizEnd

\begin{table*}
\rowcolors{6}{}{lightgray}
\renewcommand*{\arraystretch}{1.2}
\IfEbookSize{\scriptsize}{\footnotesize}
\centering
\OneColumnHSpace{-0.7in}%
\ebresizewidth{
\begin{tabular}{S[table-format = 1.0]S[table-format = 1.0]S[table-format = 9.0]
		S[table-format = 9.0]S[table-format = 2.0]S[table-format = 5.2]
		S[table-format = 4.2]S[table-format = 2.0]S[table-format = 4.2]
		S[table-format = 6.2]}
	\toprule
	\multicolumn{4}{r}{} & \multicolumn{3}{c}{\tco{-DCOLLAPSE}} &
					\multicolumn{3}{c}{\tco{-DMA=N}} \\
	\cmidrule(l){5-7} \cmidrule(l){8-10}
	\multicolumn{1}{r}{更新者} &
	    \multicolumn{1}{r}{读者} &
		\multicolumn{1}{r}{状态数} &
		    \multicolumn{1}{r}{达到的深度} &
			\multicolumn{1}{r}{\tco{-wN}} &
			    \multicolumn{1}{r}{内存 (MB)} &
				\multicolumn{1}{r}{运行时间 (s)} &
				    \multicolumn{1}{r}{\tco{N}} &
					\multicolumn{1}{r}{内存 (MB)} &
					    \multicolumn{1}{r}{运行时间 (s)} \\
	\cmidrule{1-4} \cmidrule(l){5-7} \cmidrule(l){8-10}
	1 & 1 &           376 &         95 & 12 &     0.10 & 0.00 &
		40 &    0.29 &      0.00 \\
	1 & 2 &         6 177 &        218 & 12 &     0.39 & 0.01 &
		47 &    0.59 &      0.02 \\
	1 & 3 &        99 728 &        385 & 16 &     4.60 & 0.14 &
		54 &    3.04 &      0.45 \\
        2 & 1 &        29 399 &        859 & 16 &     2.30 & 0.03 &
		55 &    0.70 &      0.13 \\
        2 & 2 &     1 071 181 &      2 352 & 20 &    49.24 & 1.45 &
		62 &    7.77 &      5.76 \\
        2 & 3 &    33 866 736 &     12 857 & 24 & 1 540.70 & 62.5 &
		69 &  111.66 &    326    \\
        3 & 1 &     2 749 453 &     53 809 & 21 &   125.25 & 4.01 &
		70 &   11.41 &     19.5  \\
        3 & 2 &   186 202 860 &    328 014 & 28 & 10 482.51 & 390 &
		77 &  222.26 &   2560    \\
	3 & 3 & 9 664 707 100 &  2 055 621 &    &          &      &
		84 & 5557.02 & 266000    \\
	\bottomrule
\end{tabular}
}
\caption{QRCU Spin 结果汇总}
\label{tab:formal:promela:QRCU Spin Result Summary}
\end{table*}

作为参考，\cref{tab:formal:promela:QRCU Spin Result Summary}
汇总了使用 \co{-DCOLLAPSE} 和 \co{-DMA=N}
编译器标志的 Spin 结果。
内存使用量是使用表中所示的最小足够搜索深度和
\co{-DMA=N} 参数获得的。
\co{-DCOLLAPSE} 运行的哈希表大小通过 \co{./pan} 的 \co{-wN}
选项进行调整，以避免在小状态空间上使用过多内存进行哈希。
因此内存使用量小于
\cref{tab:advsync:Memory Usage of QRCU Model}中所示的值，
后者的哈希表大小从默认的 \co{-w24} 开始。
运行时间来自一台 \Power{9} 服务器，这表明 \co{-DMA=N}
遭受的 CPU 开销比 \co{-DCOLLAPSE} 高约一个数量级，
但另一方面将内存开销减少了远超一个数量级。

目前为止一切顺利。
但再增加几个更新者或读者就会耗尽内存，
即使使用 \co{-DMA=N}。\footnote{
	或者，CPU 消耗会变得过大。}
那怎么办？
以下是一些可能的方法：

\begin{enumerate}
\item	看看较少数量的读者和更新者是否足以证明一般情况。
\item	手动构建正确性证明。
\item	使用更强大的工具。
\item	分而治之。
\end{enumerate}

以下各节讨论这些方法。

\subsubsection{到底需要多少读者和更新者？}
\label{sec:formal:How Many Readers and Updaters Are Really Needed?}

一种方法是仔细检查 \co{qrcu_updater()} 的 Promela 代码，
注意到唯一的全局状态变化发生在锁保护之下。
因此，在同一时刻只有一个更新者可能在修改
对读者或其他更新者可见的状态。
这意味着状态变化的任何序列都可以由单个更新者
串行执行，这是由于 Promela 进行完整的状态空间搜索。
因此，最多需要两个更新者：
一个用于修改状态，另一个用于将其搞乱。

读者的情况则不太明确，因为每个读者
仅仅执行一个单一的读端临界区然后退出。
可以论证有用的读者数量是受限的，
因为快速路径必定看到计数器中为零或者为一。
这是一个富有成效的研究方向，事实上，它导向了
下一节中描述的完整正确性证明。

\subsubsection{替代方法：
				     正确性证明}
\label{sec:formal:Alternative Approach: Proof of Correctness}

一个非正式证明~\cite{PaulMcKenney2007QRCUpatch}如下：

\begin{enumerate}
\item	如果 \co{synchronize_qrcu()} 过早退出，
	则根据定义，在 \co{synchronize_qrcu()} 的整个执行过程中
	必须至少有一个读者存在。
\item	对应该读者的计数器在此时间间隔内至少为 1。
\item	\co{synchronize_qrcu()} 代码强制
	至少一个计数器在任何时候至少为 1。
\item	以上两项意味着，如果对应该读者的计数器恰好为一，
	则另一个计数器必须大于或等于一。
	类似地，如果另一个计数器等于零，
	则对应读者的计数器必须大于或等于二。
\item	因此，在任何给定时间点，要么其中一个计数器至少为 2，
	要么两个计数器都至少为一。
\item	然而，\co{synchronize_qrcu()} 的快速路径代码
	在给定时间只能读取一个计数器。
	因此快速路径代码可能在获取第一个计数器时为零，
	但与计数器翻转竞争，使得第二个计数器被看到为一。
\item	在这样的竞态条件下，最多只能有一个读者持续存在，
	否则总和将为二或更大，
	这将导致更新者走慢速路径。
\item	但如果竞争发生在快速路径的第一次读取计数器时，
	然后在第二次读取时再次发生，
	则必须发生了两次计数器翻转。
\item	因为给定的更新者只翻转计数器一次，
	而且更新端锁阻止一对更新者并发翻转计数器，
	快速路径代码能够与翻转竞争两次的唯一方式是
	第一个更新者已经完成。
\item	但第一个更新者要等到所有先前存在的读者完成之后
	才会完成。
\item	因此，如果快速路径连续两次与计数器翻转竞争，
	所有先前存在的读者必定已经完成，
	所以走快速路径是安全的。
\end{enumerate}

当然，并非所有并行算法都有如此简单的证明。
在这种情况下，可能需要借助更强大的工具。

\subsubsection{替代方法：
				     更强大的工具}
\label{sec:formal:Alternative Approach: More Capable Tools}

虽然 Promela 和 Spin 十分有用，
但有更强大的工具可用，特别是用于验证硬件的工具。
这意味着如果能将你的算法翻译为硬件设计的 VHDL 语言——
对于非常底层的并行算法通常如此——
则可以将这些工具应用于你的代码
（例如，这曾被用于第一个实时 RCU 算法）。
然而，这些工具可能相当昂贵。

虽然商业化多处理器的出现可能最终会催生
具有强大状态空间缩减能力的自由软件模型检查器，
但这目前还没有什么大的帮助。

另外，Spin 有一些特性支持需要固定内存量的近似搜索，
然而，在验证并行算法时，我始终无法让自己信任这种近似方法。

另一种方法可能是分而治之。

\subsubsection{替代方法：
				     分而治之}
\label{sec:formal:Alternative Approach: Divide and Conquer}

通常可以将较大的并行算法分割成较小的片段，
然后分别证明它们。
例如，一个 100 亿状态的模型可能被分解为
一对 100,000 状态的模型。
采用这种方法不仅使 Promela 等工具更容易验证你的算法，
也会使你的算法更容易理解。

\subsubsection{QRCU 真的正确吗？}
\label{sec:formal:Is QRCU Really Correct?}

QRCU 真的正确吗？
我们有一个基于 Promela 的机械证明和一个手工证明，
二者都说它是正确的。
然而，\pplsur{Jade}{Alglave} 等人的一篇论文~\cite{JadeAlglave2013-cav}
却持相反意见（见论文第 12 页底部的 Section~5.1）。
到底怎么回事？

事实证明两者都是正确的！
当 QRCU 被添加到形式验证基准测试套件中时，
它的内存屏障被遗漏了，从而产生了一个有缺陷的 QRCU 版本。
所以这里真正的新闻是，许多形式验证工具
错误地证明了这个有缺陷的 QRCU 是正确的。
这就是为什么形式验证工具本身应该使用
被验证代码的注入缺陷版本来进行测试。
如果给定的工具无法找到注入的缺陷，那么它显然是不可信的。

\QuickQuiz{
	但不同的形式验证工具通常被设计为定位特定类别的缺陷。
	例如，很少有形式验证工具能发现规范中的错误。
	所以这个``显然不可信''的判断是不是有点苛刻？
}\QuickQuizAnswer{
	许多形式验证工具确实在某些方面是专门化的。
	例如，Promela 不处理现实的内存模型（memory model）
	（尽管可以将它们编程到
	\IX{Promela}~\cite{Desnoyers:2013:MSM:2506164.2506174} 中），
	\IXacr{cbmc}~\cite{EdmundClarke2004CBMC} 不检测概率性
	挂起和死锁（deadlock），
	\IX{Nidhugg}~\cite{CarlLeonardsson2014Nidhugg} 不检测
	涉及数据非确定性的缺陷。
	但这意味着这些工具不能被信任来发现
	它们不是为定位而设计的缺陷。

	因此，创建形式验证工具的人应该
	``在标签上说实话''，明确指出他们的工具
	能够和不能检测到哪些类别的缺陷。
	否则，当一个从业者第一次发现某个工具
	未能检测到一个缺陷时，该从业者很可能会
	对该工具进行极其严厉和极其公开的谴责。
	是的，是的，展示最好的一面确实有道理，
	但在没有适当免责声明的情况下走得太远，
	很容易触发负面反应的地雷，
	你的工具可能能也可能不能从中恢复。

	你已经被警告了！
}\QuickQuizEnd

因此，如果你确实打算使用 QRCU，请务必小心。
它的正确性证明本身可能正确也可能不正确。
这是形式验证不太可能完全取代测试的原因之一，
正如 Donald Knuth 很早以前就指出的那样。

\QuickQuiz{
	鉴于我们对此处描述的 QRCU 算法有两个独立的正确性证明，
	而且不正确性证明涵盖的是已知为不同的算法，
	为什么还有怀疑的余地？
}\QuickQuizAnswer{
	总是有怀疑的余地。
	在这种情况下，重要的是要记住这两个正确性证明
	先于现实世界内存模型的形式化，
	这就提出了这两个证明可能基于
	不正确的内存排序假设的可能性。
	此外，由于两个证明都是由同一个人构建的，
	它们很可能包含一个共同的错误。
	再说一次，总是有怀疑的余地。
}\QuickQuizEnd
